\section{Methodology}
\label{sec:method}

We formalize the adaptive question selection problem and describe the 14 algorithms evaluated in our study.

\subsection{Problem Formulation}
\label{sec:formulation}

\paragraph{Student model.}
Each student maintains a knowledge vector $\mathbf{k} = [k_1, \ldots, k_K] \in [0,1]^K$ over $K$ topics.
When quizzed on topic $i$, the student answers correctly with probability $k_i$ (Bernoulli trial).
Knowledge evolves according to two dynamics:

\noindent\textit{Learning.}
After answering a question on topic $i$:
\begin{equation}
    k_i \leftarrow
    \begin{cases}
        k_i + \alpha \cdot (1 - k_i), & \text{if correct}, \\
        k_i - \beta \cdot k_i, & \text{if incorrect},
    \end{cases}
    \label{eq:learning}
\end{equation}
where $\alpha = 0.12$ is the learning rate and $\beta = 0.02$ is the incorrect-answer penalty.
The asymptotic form ensures diminishing returns as knowledge approaches mastery.

\noindent\textit{Forgetting.}
At each timestep, all topics \emph{not} quizzed undergo Ebbinghaus-inspired exponential decay~\citep{ebbinghaus1885memory}:
\begin{equation}
    k_j \leftarrow k_{\mathrm{base}} + (k_j - k_{\mathrm{base}}) \cdot e^{-\lambda}, \quad \forall j \neq i,
    \label{eq:forgetting}
\end{equation}
where $k_{\mathrm{base}} = 0.10$ is the retention floor and $\lambda > 0$ is the per-step decay rate.
Higher $\lambda$ corresponds to faster forgetting.

\paragraph{RMAB formulation.}
We cast this problem as a restless multi-armed bandit~\citep{whittle1988restless}.
Each of the $K$ topics is an arm.
The state of arm~$i$ at time~$t$ is its knowledge level $k_i^{(t)} \in [0,1]$.
At each timestep the learner selects one arm to activate (quiz).
The activated arm transitions according to \eqref{eq:learning}, while all passive arms decay according to \eqref{eq:forgetting}.
The objective is to maximize cumulative mean knowledge:
\begin{equation}
    \max_{\pi} \;\; \E\!\left[\frac{1}{T} \sum_{t=1}^{T} \bar{k}^{(t)}\right], \quad \text{where } \bar{k}^{(t)} = \frac{1}{K}\sum_{i=1}^{K} k_i^{(t)}.
    \label{eq:objective}
\end{equation}
This RMAB is intractable in general (PSPACE-hard~\citep{papadimitriou1999complexity}), motivating the heuristic algorithms below.

\subsection{Baseline Algorithms}
\label{sec:baselines}

We evaluate seven baseline algorithms spanning classical bandits, heuristics, and a myopic greedy policy.

\paragraph{Random.}
Selects a topic uniformly at random at each step.
Serves as the primary baseline; all improvements are reported relative to Random.

\paragraph{UCB1~\citep{auer2002finite}.}
Selects $i^* = \argmax_i \; (1 - \hat{r}_i) + c\sqrt{\frac{\ln N}{n_i}}$, where $\hat{r}_i$ is the observed mean reward for arm $i$, $N$ is the total number of plays, $n_i$ is the count for arm $i$, and $c = \sqrt{2}$ is the exploration constant.
The inverted reward $(1 - \hat{r}_i)$ targets weak topics.

\paragraph{Thompson Sampling~\citep{thompson1933likelihood}.}
Maintains a Beta posterior $\text{Beta}(\alpha_i, \beta_i)$ for each arm.
At each step, samples $\theta_i \sim \text{Beta}(\alpha_i, \beta_i)$ and selects $i^* = \argmax_i (1 - \theta_i)$.

\paragraph{$\varepsilon$-Greedy.}
With probability $\varepsilon = 0.1$ selects a random arm; otherwise selects the arm with the lowest observed mean reward.

\paragraph{Sliding-Window UCB (SW-UCB)~\citep{garivier2011upper}.}
Restricts reward estimates to the most recent $\tau$ observations per arm, adapting to non-stationarity.
Uses $\tau = 50$ in our experiments.

\paragraph{Leitner System~\citep{leitner1972so}.}
Implements the classical box-based spaced repetition heuristic.
Topics start in box~1 and advance on correct answers; incorrect answers reset to box~1.
Lower boxes are reviewed more frequently.

\paragraph{Greedy-Min.\footnote{The Greedy-Min policy selects the category with lowest current knowledge. Despite having access to the true knowledge state, it is myopic---it does not plan ahead or account for future forgetting dynamics.}}
Has access to the true knowledge vector $\mathbf{k}^{(t)}$ at each step and selects $i^* = \argmin_i \; k_i^{(t)}$---the topic with the lowest current knowledge.
Despite perfect information, this myopic greedy policy does not account for future forgetting.

\subsection{Novel Algorithms}
\label{sec:novel}

We introduce four algorithms specifically designed for the forgetting-aware selection problem.

\paragraph{Forgetting-UCB (F-UCB).}
F-UCB augments the standard UCB score with a decay-aware urgency term.
Let $t_i$ be the number of steps since arm $i$ was last played and $\hat{r}_i$ the observed mean reward.
The score is:
\begin{equation}
    s_i^{\text{F-UCB}} = \underbrace{(1 - \hat{r}_i) \cdot e^{-\lambda t_i}}_{\text{weighted weakness}} + \underbrace{\gamma \cdot (1 - e^{-\lambda t_i})}_{\text{decay urgency}} + \underbrace{c\sqrt{\frac{\ln N}{n_i}}}_{\text{exploration}},
    \label{eq:fucb}
\end{equation}
where $\gamma = 1.0$ controls the urgency weight and $\lambda$ is the known decay rate.
The first term downweights the observed weakness estimate as it becomes stale; the second term increases urgency for arms that have been decaying longer; the third is the standard UCB exploration bonus.

\paragraph{BKT-Bandit.}
BKT-Bandit maintains a Bayesian posterior over each arm's mastery probability, augmented with forgetting.
Each arm $i$ has parameters $(\alpha_i, \beta_i)$ of a Beta distribution, updated on observations:
\begin{equation}
    \alpha_i \leftarrow \alpha_i + \mathbb{1}[\text{correct}], \quad
    \beta_i \leftarrow \beta_i + \mathbb{1}[\text{incorrect}].
\end{equation}
Between observations, the posterior is decayed toward uncertainty:
\begin{equation}
    \alpha_i \leftarrow 1 + (\alpha_i - 1) \cdot e^{-\lambda}, \quad
    \beta_i \leftarrow 1 + (\beta_i - 1) \cdot e^{-\lambda}.
\end{equation}
The exponential posterior shrinkage toward the prior can be interpreted as a discounted posterior~\citep{raj2017taming}, where the effective sample size of past observations decays geometrically with rate $e^{-\lambda}$. This is equivalent to a Bayesian update under a hazard model where the student's knowledge state may change (due to forgetting) with probability $1 - e^{-\lambda}$ per timestep, justifying the use of a time-varying posterior.

The selection score combines weakness with uncertainty:
\begin{equation}
    s_i^{\text{BKT}} = \underbrace{(1 - \mu_i)}_{\text{weakness}} + \underbrace{c \cdot \sigma_i}_{\text{uncertainty bonus}},
    \label{eq:bkt}
\end{equation}
where $\mu_i = \frac{\alpha_i}{\alpha_i + \beta_i}$ and $\sigma_i = \sqrt{\frac{\alpha_i \beta_i}{(\alpha_i + \beta_i)^2(\alpha_i + \beta_i + 1)}}$ are the posterior mean and standard deviation.

\paragraph{Whittle Advantage.}
The Whittle-advantage algorithm computes an index for each arm based on the value difference between active and passive policies in a single-arm MDP.
For arm $i$ with current state $k_i$, the Whittle index approximates:
\begin{equation}
    W_i = V_i^{\text{active}}(k_i) - V_i^{\text{passive}}(k_i),
    \label{eq:whittle}
\end{equation}
where $V^{\text{active}}$ is the value of quizzing the arm (applying \eqref{eq:learning}) and $V^{\text{passive}}$ is the value of leaving it idle (applying \eqref{eq:forgetting}).
We solve the single-arm MDP via value iteration with a discretized state space.
The arm with the highest Whittle index is selected at each step.
We evaluate two variants: \textbf{PD-Whittle} (probability-dependent) and \textbf{Adapt-Whittle} (adaptive estimation).

\paragraph{MetaSelector.}
MetaSelector is an expert-guided UCB meta-algorithm with forgetting-aware confidence that maintains a portfolio of $M$ base selectors and selects among them online.
Each base selector $j \in \{1, \ldots, M\}$ is one of the algorithms above.
MetaSelector maintains counts $N_j$ and mean rewards $\bar{r}_j$ for each expert and selects:
\begin{equation}
    j^* = \argmax_j \; \bar{r}_j + c_{\text{meta}}\sqrt{\frac{\ln N}{N_j}},
    \label{eq:meta}
\end{equation}
where the reward signal is the student's average knowledge after acting on expert $j$'s recommendation.
The selected expert's arm recommendation is then executed.
MetaSelector achieves principled regime interpolation: it automatically shifts weight toward whichever base algorithm performs best in the current $(K, \lambda)$ regime, without requiring prior knowledge of which algorithm to use.
We explored formal expert aggregation approaches (EXP4, Corral~\citep{agarwal2017corralling}) but found that the partial-observability reward signal---binary correctness rather than knowledge gain---prevented reliable expert tracking.
Its unified scoring combines weakness, decay, and uncertainty:
\begin{equation}
    s_i^{\text{unified}} = (1 - \mu_i) \cdot e^{-\lambda t_i} + 0.5 \cdot (1 - e^{-\lambda t_i}) + c \cdot \sigma_i.
    \label{eq:unified}
\end{equation}
